\uuid{Mpau}
\exo7id{118}
\titre{exo7 118}
\auteur{ridde}
\organisation{exo7}
\datecreate{1999-11-01}
\isIndication{false}
\isCorrection{false}
\chapitre{Logique, ensemble, raisonnement}
\sousChapitre{Logique}
\cpge{mpsi-mp2i.raisonnement-vocabulaire-ensembliste.a, pcsi.raisonnement-vocabulaire-ensembliste.a, ptsi.raisonnement-vocabulaire-ensembliste.a, mpsi-mp2i.raisonnement-vocabulaire-ensembliste.b, pcsi.raisonnement-vocabulaire-ensembliste.b, ptsi.raisonnement-vocabulaire-ensembliste.b, bcpst1.logique-ensembles-raisonnement.a, bcpst1.logique-ensembles-raisonnement.b}
\module{Algèbre}
\niveau{L1}
\difficulte{}

\contenu{
\question{
Si $P (x)$ est une proposition d\'ependant de $x \in X$, on note $\overline{P} =
\left\{ x \in X / P (x) \text{ est vraie}\right\}$.
Exprimer en fonction de $\overline{P}$ et $\overline{Q}$ les ensembles $\overline{\neg P},
 \overline{P \wedge Q}, \overline{P \vee Q}, \overline{P \Rightarrow Q},
\overline{P \Leftrightarrow Q}$.
}
}
